% Copyright (c) 2026 Benjamin Stanley Frohman.
% This paper is licensed under the Creative Commons Attribution 4.0
% International License (CC-BY-4.0). The Lean code it describes is Apache-2.0.
\documentclass[11pt]{article}
\usepackage[utf8]{inputenc}
\usepackage[T1]{fontenc}
\usepackage{lmodern}
\usepackage{underscore}
\usepackage{newunicodechar}
\usepackage{amsmath,amssymb,amsthm}
\usepackage[margin=1in]{geometry}
\usepackage{booktabs}
\usepackage{fancyvrb}
\usepackage[hidelinks]{hyperref}
\newunicodechar{ω}{\ensuremath{\omega}}
\newunicodechar{κ}{\ensuremath{\kappa}}
\newunicodechar{δ}{\ensuremath{\delta}}
\newunicodechar{ℝ}{\ensuremath{\mathbb{R}}}
\newunicodechar{∀}{\ensuremath{\forall}}
\newunicodechar{∃}{\ensuremath{\exists}}
\newunicodechar{¬}{\ensuremath{\neg}}
\newunicodechar{↔}{\ensuremath{\leftrightarrow}}
\newunicodechar{→}{\ensuremath{\to}}
\newunicodechar{∧}{\ensuremath{\wedge}}
\newunicodechar{∨}{\ensuremath{\vee}}
\newunicodechar{₁}{\ensuremath{_1}}
\newunicodechar{₂}{\ensuremath{_2}}
\newunicodechar{₃}{\ensuremath{_3}}

\newtheorem{theorem}{Theorem}
\newtheorem{proposition}[theorem]{Proposition}
\theoremstyle{definition}
\newtheorem{definition}[theorem]{Definition}
\theoremstyle{remark}
\newtheorem{remark}[theorem]{Remark}

\newcommand{\lean}[1]{\texttt{\upshape #1}}
\fvset{fontshape=n,fontseries=m}
\newcommand{\R}{\mathbb{R}}
\newcommand{\TF}{\mathrm{TF}}
\newcommand{\KFT}{\mathcal{K}_{\mathrm{FT}}}

\title{The Frohmanian Tether kernels in Lean:\\
non-antisymmetry, energy-direction degeneracy,\\
and the Beltrami alignment zero set}
\author{Benjamin Stanley Frohman\\
\small Independent Researcher, Texas, United States\\
\small ORCID \href{https://orcid.org/0009-0006-7068-3718}{0009-0006-7068-3718}\quad
\href{mailto:frohmanbenjamin@gmail.com}{frohmanbenjamin@gmail.com}}
\date{November 23, 2026}

\begin{document}
\maketitle

\begin{abstract}
We record, with machine-checked Lean~4/Mathlib proofs, exact pointwise
algebraic facts about the densities of the ``Frohmanian Symplectic Tether''
($\TF$) kernels written in the author's May~2026 notes and the October~2026
kernel $\KFT$. (i) The displayed associated 2-form density
$\omega\cdot(\delta u\times\delta v)-\kappa|\omega|^2(\delta u\cdot\delta v)$ is
not antisymmetric, so the claim that it is a bilinear antisymmetric correction
to Arnold's Lie--Poisson bracket is false. (ii) In the projected-kernel reading,
the tether term along the kinetic-energy direction $\delta H=u$ is identically
zero, so it adds nothing to the Navier--Stokes equations. (iii) The October
kernel satisfies $\KFT(\omega;u,u)=|\omega\times u|^2$, which vanishes exactly
when $u\parallel\omega$ (Beltrami, i.e.\ velocity--vorticity, alignment) and is
not identically zero. (iv) $|\sin\theta|$ and $|\tau-\eta|$ are different
quantities for vorticity directions $\tau,\eta$. \emph{This is not a regularity
result}: no bound, blow-up criterion, or regularity statement for Navier--Stokes
solutions is proved.
\end{abstract}

\section{Notation and scope}
We use textbook notation: $u$ is the velocity, $\omega=\nabla\times u$ the
vorticity, $\nu$ the kinematic viscosity, $\|\omega\|_{L^2}^2$ the enstrophy,
$H=\tfrac12\int|u|^2\,dx$ the kinetic energy (so its variational derivative is
$\delta H=u$), $\tau=\omega/|\omega|$ the vorticity direction, and $\theta$ the
angle between two vorticity vectors. The incompressible Navier--Stokes
equations are $\partial_t u+(u\cdot\nabla)u=-\nabla p+\nu\Delta u$,
$\nabla\cdot u=0$.

Every Lean statement below concerns vectors in $\R^3$ (Lean type
\lean{Fin 3 → ℝ}) at a single point; $a\cdot b$ is \lean{dot a b}, $a\times b$ is
\lean{cross a b}, and $e_1,e_2,e_3$ are the standard basis vectors
\lean{e₁}, \lean{e₂}, \lean{e₃}. No integral over $\R^3$ or $\mathbb{T}^3$
is formalized. A pointwise density that fails a property (for example
antisymmetry) at one point cannot acquire it after integration against every
test field, which is why the pointwise statements suffice for the refutation in
Section~\ref{sec:antisym}.

\section{Definitions (Lean: \lean{TFFrohmanianTether/Basic.lean})}
\begin{definition}[\lean{TF}]
For $\kappa\in\R$ and $\omega,\delta u,\delta v\in\R^3$,
\[\TF_\kappa(\omega;\delta u,\delta v)=\underbrace{\omega\cdot(\delta u\times\delta v)}_{\text{\lean{ω_KKS}}}
-\kappa\underbrace{|\omega|^2(\delta u\cdot\delta v)}_{\text{\lean{omega_sq_dot}}}.\]
\end{definition}
\begin{definition}[\lean{Pi_u}, \lean{P_δC_perp}]
$\Pi_u v=v-\dfrac{v\cdot u}{u\cdot u}\,u$ if $u\neq0$ and $\Pi_0v=v$.
$P_\omega^\perp v=\Pi_\omega v$ is the projection off the Casimir direction
$\omega$.
\end{definition}
\begin{definition}[\lean{K_FT}]
$\KFT(\omega;v,w)=|\omega|^2\,\big(P_\omega^\perp v\big)\cdot\big(P_\omega^\perp w\big)$.
\end{definition}

\section{The displayed 2-form is not antisymmetric}\label{sec:antisym}
The May~21, 2026 note (\texttt{vFinal.txt}, sha256
\texttt{d5b6508e\ldots c6e05}) states: ``We derive the Frohmanian Symplectic
Tether TF as the canonical minimal bilinear antisymmetric correction to Arnold's
Lie--Poisson bracket'' and ``The associated 2-form is
$\TF(\delta u,\delta v)=\int_{\mathbb{T}^3}\omega\cdot(\delta u\times\delta v)\,d\lambda
-\kappa\int_{\mathbb{T}^3}|\omega|^2(\delta u\cdot\delta v)\,d\lambda$.''
The same formula appears in the May~14 paper
(sha256 \texttt{98717d94\ldots 026ff}).

\begin{definition}[\lean{paper_claim}, \lean{TFFrohmanianTether/Semantics.lean}]
The literal transcription of the antisymmetry claim is
\begin{Verbatim}[fontsize=\small]
def paper_claim : Prop :=
  ∀ (κ : ℝ) (ω δu δv : Fin 3 → ℝ),
    TF κ ω δu δv = - TF κ ω δv δu
\end{Verbatim}
\end{definition}

\begin{theorem}[\lean{TF_not_antisymmetric}, \lean{paper_claim_refuted}]
It is not true that $\TF_\kappa(\omega;\delta u,\delta v)=-\TF_\kappa(\omega;\delta v,\delta u)$
for all $\kappa,\omega,\delta u,\delta v$. Witness: $\kappa=1$,
$\omega=\delta u=\delta v=e_1$ gives $\TF=-1$ but $-\TF=1$.
\begin{Verbatim}[fontsize=\small]
theorem TF_not_antisymmetric :
    ¬ (∀ (κ : ℝ) (ω δu δv : Fin 3 → ℝ),
        TF κ ω δu δv = - TF κ ω δv δu)
theorem paper_claim_refuted : ¬ paper_claim
\end{Verbatim}
\end{theorem}

\begin{theorem}[\lean{TF_not_symmetric}]
Nor is it symmetric: $\kappa=0$, $\omega=e_1$, $\delta u=e_2$, $\delta v=e_3$
gives $\TF(\delta u,\delta v)=1$ and $\TF(\delta v,\delta u)=-1$.
\end{theorem}

\begin{proposition}[\lean{intended_correspondence}]
With \lean{formal_claim := False}, \lean{formal_claim ↔ paper_claim}. The
proof uses the witness of \lean{TF_not_antisymmetric}; it is not
\lean{Iff.rfl}.
\end{proposition}

The supporting lemmas \lean{ω_KKS_antisymm}
($\omega\cdot(\delta u\times\delta v)=-\omega\cdot(\delta v\times\delta u)$),
\lean{omega_sq_dot_symm} (the metric term is symmetric), \lean{TF_split},
\lean{cross_antisymm}, \lean{dot_comm}, \lean{dot_neg_right},
\lean{dot_smul_right} and \lean{eq_zero_of_dot_self_eq_zero} ($a\cdot a=0\Rightarrow a=0$)
are also proved. The failure is exactly the symmetric metric term
$-\kappa|\omega|^2(\delta u\cdot\delta v)$, which is not alternating.

\section{The projected tether adds nothing along the energy direction}
\begin{theorem}[\lean{Pi_u_self}, \lean{Pi_u_pairing_with_u}]
$\Pi_u u=0$, hence $(\Pi_u\,\delta F)\cdot(\Pi_u u)=0$ for all $u,\delta F$.
\end{theorem}
\begin{theorem}[\lean{projected_TF_adds_nothing_along_H}, \lean{Pi_u_tether_contribution_along_H}]
For all $\kappa,\omega,u,\delta F$:
$-\kappa\,|\omega|^2\,(\Pi_u\,\delta F)\cdot(\Pi_u u)=0$.
\begin{Verbatim}[fontsize=\small]
theorem projected_TF_adds_nothing_along_H (κ : ℝ) (ω u δF : Fin 3 → ℝ) :
    -κ * omega_sq_dot ω (Pi_u u δF) (Pi_u u u) = 0
\end{Verbatim}
\end{theorem}
\begin{remark}
With $\delta H=u$, the tether contribution to $dF/dt=\{F,H\}$ in the
projected-kernel reading is identically zero, so it adds no term to the
Navier--Stokes vorticity equation from which a bound could be derived. The
lemmas \lean{omega_sq_dot_Pi_u_self} and \lean{omega_sq_dot_Pi_u_self_weight}
state the same vanishing with weight $u$ and with an independent weight
$\omega$. Without the projection the term need not vanish:
\lean{omega_sq_dot_not_degenerate_on_u} gives $\omega,u$ with
$|\omega|^2(u\cdot u)\neq0$.
\end{remark}

\section{The October kernel is the Beltrami alignment defect}
\begin{theorem}[\lean{P_δC_perp_self}, \lean{K_FT_vanishes_on_δC}, \lean{K_FT_vanishes_on_δC_right}]
$P_\omega^\perp\omega=0$, so $\KFT(\omega;\omega,v)=\KFT(\omega;v,\omega)=0$.
\end{theorem}
\begin{theorem}[\lean{K_FT_uu_eq_cross_sq}]
For all $\omega,u\in\R^3$, $\KFT(\omega;u,u)=|\omega\times u|^2$.
\begin{Verbatim}[fontsize=\small]
theorem K_FT_uu_eq_cross_sq (ω u : Fin 3 → ℝ) :
    K_FT ω u u = dot (cross ω u) (cross ω u)
\end{Verbatim}
\end{theorem}
\begin{proof}[Proof idea]
If $\omega=0$ both sides vanish. Otherwise
$|\omega|^2|P_\omega^\perp u|^2=|\omega|^2|u|^2-(u\cdot\omega)^2=|\omega\times u|^2$
(Lagrange's identity), checked in Lean by \lean{field_simp} and \lean{ring}.
\end{proof}
\begin{theorem}[\lean{K_FT_uu_eq_zero_of_collinear}, \lean{K_FT_uu_eq_zero_iff}]
$\KFT(\omega;u,u)=0$ if and only if $\omega=0$ or $u=c\,\omega$ for some
$c\in\R$, i.e.\ exactly at velocity--vorticity (Beltrami) alignment
$u\parallel\omega$.
\begin{Verbatim}[fontsize=\small]
theorem K_FT_uu_eq_zero_of_collinear (ω u : Fin 3 → ℝ) (c : ℝ)
    (hu : u = fun i => c * ω i) : K_FT ω u u = 0
theorem K_FT_uu_eq_zero_iff (ω u : Fin 3 → ℝ) :
    K_FT ω u u = 0 ↔ (ω = fun _ => 0) ∨ ∃ c : ℝ, u = fun i => c * ω i
\end{Verbatim}
\end{theorem}
\begin{theorem}[\lean{K_FT_uu_not_identically_zero}, \lean{K_FT_uu_nonneg}]
There are $\omega,u$ with $\KFT(\omega;u,u)\neq0$ (take $\omega=e_1$, $u=e_2$),
and $\KFT(\omega;u,u)\ge0$ always.
\end{theorem}
\begin{remark}
Integrated over a domain, $\int|\omega\times u|^2\,dx$ is therefore a measure of
how far a flow is from Beltrami alignment. That integrated equivalence is
\emph{not} formalized here, and nothing here shows that Navier--Stokes
solutions become or stay aligned.
\end{remark}

\section{Vorticity direction: $|\sin\theta|$ versus $|\tau-\eta|$}
Constantin and Fefferman~\cite{CF93} proved regularity under a Lipschitz
condition on the vorticity direction where $|\omega|$ is large, and
Beirão da Veiga and Berselli~\cite{BVB02} showed that the Hölder-$\tfrac12$
condition $|\sin\theta(x,y)|\le C|x-y|^{1/2}$ suffices. Here
$|\sin\theta|=|\tau(x)\times\tau(y)|$ for unit $\tau$.
\begin{theorem}[\lean{sin_theta_sq_of_parallel}, \lean{sin_theta_sq_of_antiparallel}]
$|\tau\times\tau|^2=0$ and $|\tau\times(-\tau)|^2=0$ for every $\tau\in\R^3$.
\end{theorem}
\begin{theorem}[\lean{antiparallel_direction_gap}]
For $\tau=e_1$ and $\eta=-e_1$: $|\tau\times\eta|^2=0$, so $\sin\theta=0$,
while $|\tau-\eta|^2=4$.
\begin{Verbatim}[fontsize=\small]
theorem antiparallel_direction_gap :
    sin_theta_sq e₁ (fun i => -e₁ i) = 0 ∧
      dot (fun i => e₁ i - (-e₁ i)) (fun i => e₁ i - (-e₁ i)) = 4
\end{Verbatim}
\end{theorem}
So a bound on $|\sin\theta|$ and a bound on $|\tau(x)-\tau(y)|$ are different
conditions and cannot be substituted for one another.

\section{What is not proved}
No statement here is a regularity theorem. In particular, neither the
Constantin--Fefferman / Beirão da Veiga--Berselli direction condition nor
velocity--vorticity alignment is proved for every smooth divergence-free
initial datum, and parts (A) and (B) of the Clay problem
statement~\cite{Fef06} are not addressed.

\section{Reproducibility and certificate}\label{sec:cert}
Source: \url{https://github.com/BenFrohman/tf-frohmanian-tether-audit}.
Pins:
\begin{Verbatim}[fontsize=\small]
lean-toolchain      leanprover/lean4:v4.30.0-rc1
Mathlib rev         0c154d67103f74be3a0f2c509f72ccbf5be9f2a7
lake-manifest.json  sha256 b81d814b7d2d2d2821b03ed679324d5e
                           7ec7da674e0a977352d01c88cbe8ccb8
\end{Verbatim}
Build:
\begin{Verbatim}[fontsize=\small]
lake exe cache get
lake build
lake env lean TFFrohmanianTether/AxiomAudit.lean
\end{Verbatim}
The official certificate is the GitHub Actions run of
\texttt{.github/workflows/lean-kernel.yml} on the exact commit and its artifact
\texttt{lean-kernel-certificate}. A pass means \texttt{lake build} succeeds,
there are 0 \texttt{sorry}, and every \texttt{\#print axioms} lists only
\texttt{propext}, \texttt{Classical.choice}, \texttt{Quot.sound}.

%CERT-BEGIN
\noindent\textbf{CI certificate run:} to be filled in from the first green CI run.
%CERT-END

\section*{License}
This paper: CC-BY-4.0. Lean code: Apache-2.0. Copyright \copyright\ 2026
Benjamin Stanley Frohman.

\begin{thebibliography}{9}
\bibitem{Arn66} V.~I.~Arnold, Sur la géométrie différentielle des groupes de
Lie de dimension infinie et ses applications à l'hydrodynamique des fluides
parfaits, \emph{Ann. Inst. Fourier} 16 (1966), 319--361.
\bibitem{CF93} P.~Constantin, C.~Fefferman, Direction of vorticity and the
problem of global regularity for the Navier--Stokes equations,
\emph{Indiana Univ. Math. J.} 42(3) (1993), 775--789,
doi:10.1512/iumj.1993.42.42034.
\bibitem{BVB02} H.~Beirão da Veiga, L.~C.~Berselli, On the regularizing effect
of the vorticity direction in incompressible viscous flows,
\emph{Differential Integral Equations} 15(3) (2002), 345--356,
doi:10.57262/die/1356060864.
\bibitem{Fef06} C.~L.~Fefferman, Existence and smoothness of the
Navier--Stokes equation, in \emph{The Millennium Prize Problems}, Clay Math.
Inst., 2006, 57--67.
\bibitem{Mathlib} The mathlib Community, The Lean mathematical library,
\emph{CPP 2020}, 367--381.
\end{thebibliography}
\end{document}
